\subsection{Exact corrected Hodge flags at the correction primes}
\label{sec:geometry-exact-corrected-flags}

We retain the original full-level model, cubical line, torsion
characters and locally free modules of Section~\ref{sec:geometry-hodge-specialization}.
At a correction prime \(p\), write \(R=R_p\ge4\) and \(m=R-2\).
The following argument uses all torsion levels only over the injective
valued test fields. It does not enlarge the prescribed level or the
function field of the global model.

\begin{lemma}[A compatible polarized test tower]
\label{lem:corrected-ht-test-tower}
At an injective height-one test point above \(p\), let \(C\) be the
complete algebraically closed valued field, \(O_C\) its valuation ring,
and \(G\) the pulled-back semiabelian family of dimension \(g\) with
abelian generic fibre. There is a \(p\)-divisible group \(B/O_C\) of
height \(2g\) such that:
\begin{enumerate}
\item its generic torsion at every level is the original generic
torsion \(G_C[p^n]\), with all transition maps and coefficient actions;
\item its invariant cotangent module is \(\omega_G\), with the
identification at level \(n\) given by (HS5);
\item the original alternating Weil pairings extend compatibly to
a morphism \(\lambda:B\to B^D\) whose Tate-module map is an
isomorphism preserving the original pairing.
\end{enumerate}
The kernel of the characteristic-zero Hodge--Tate map
\[
 H=\alpha_{B^D}\circ T_p(\lambda):
 T_p(B)\otimes_{\mathbb Z_p}C\longrightarrow\omega_G\otimes_{O_C} C
\]
is Lagrangian for the original pairing (using a compatible system of
roots of unity to trivialize the Tate twist).
\end{lemma}

\begin{proof}
Work first over the henselization of the height-one discrete valuation
ring of the original model. For every \(n\ge1\), the finite-part theorem
used in (HS3)--(HS4) gives the finite flat open and closed subgroup
\(\mathcal H_n\) of \(G[p^n]\) containing the entire closed fibre.
Pull these subgroups back to \(O_C\). Their complements still have
empty closed fibre, and the finite-part construction is functorial:
the inclusion and multiplication-by-\(p\) maps preserve the finite
parts. Indeed their restrictions from a finite scheme cannot meet
the complementary open and closed component if they do not meet
it on the closed fibre. The inclusion
\(\mathcal H_n\hookrightarrow\mathcal H_{n+1}\) is closed.

Multiplication by \(p\) maps \(\mathcal H_{n+1}\) faithfully flatly
onto \(\mathcal H_n\). To justify flatness, on the special fibre
these are the full semiabelian torsion schemes. The map there is
finite faithfully flat of rank \(p^{2g-r}\), where \(r\) is the
torus rank of that fibre. Over the henselian DVR both coordinate
algebras are finite free, of ranks \(p^{(n+1)(2g-r)}\) and
\(p^{n(2g-r)}\). The finite special-fibre algebra of the target is
Artinian, and the source is free of the displayed relative rank
over it. Lift a basis over that algebra. Nakayama's lemma gives
a surjection from the free module of the same relative rank over
the integral target coordinate algebra. The two underlying free
DVR modules have equal rank, so the surjection is an isomorphism.
This proves faithful flatness before base change, and hence after it.
In particular the generic finite parts map onto each other.

For each \(n\) apply the explicit induced-coset construction (HS4)
to \(T_n=G_C[p^n](C)\) and \(H_n=\mathcal H_n(C)\), obtaining \(B_n\).
An integral section of \(G[p^n]\) factors through \(\mathcal H_n\):
its closed point belongs to the latter, and its inverse image of
the complementary open and closed part is empty. Consequently
\(H_{n+1}\cap T_n=H_n\). The inclusions of generic labels therefore
induce closed immersions \(B_n\hookrightarrow B_{n+1}\), using the
finite-part inclusions on the individual copies. The maps
\(p:T_{n+1}\twoheadrightarrow T_n\) induce maps
\(B_{n+1}\twoheadrightarrow B_n\), surjective on the coset sets
and on each copy finite faithfully flat, by the preceding paragraph.
The choices of representatives cancel by the translation formula
(HS4). In particular these maps satisfy all their compatibility
relations; alternatively check the relations on the generic fibres
and use flatness and separatedness.

The kernel of \(p^n:B_{n+1}\to B_1\) is finite locally free of
rank \(p^{2gn}\). The inclusion of \(B_n\) factors through it and
is a closed immersion of finite locally free schemes of equal rank,
hence an isomorphism. This proves the Barsotti--Tate axioms for
the system \(B_n\); the same argument gives the usual exact sequences
at arbitrary levels. It is important here that the ranks of
\(B_n\) are those of the full generic abelian torsion, not those
of the intrinsic special-fibre semiabelian torsion.
The identity copy has cotangent \(\omega_G/p^n\omega_G\).
These identifications commute with the cotangent restrictions, so
their limit identifies \(\omega_B\) with \(\omega_G\).
In particular the dimension of \(B\) is \(g\). Its dual has dimension
\(2g-g=g\).

At every level the original cubical pairing extends by (HS6).
The resulting maps \(B_n\to B_n^D\) are compatible: compatibility
of the Weil pairings with inclusion and multiplication by \(p\)
holds on the generic fibres, and each equation extends between the
finite flat schemes. This yields \(\lambda\). The map of Tate modules
is an isomorphism by the original generic principal polarization;
we have not asserted that each finite-level \(\lambda_n\) is a
scheme-theoretic isomorphism. The coefficient action extends by
the same finite-part and coset construction.

Apply the integral Hodge--Tate complex theorem to \(B\) and \(B^D\)
(Fargues, \emph{Ann. Sci. ENS} 44 (2011), \S5.3.2, Theorem 2,
pp.~934--935; see also the finite-group theorem in \S5.4
already used in (HS8)). Over \(C\) it gives the exact sequence
\[
 0\longrightarrow(\omega_B\otimes C)^\vee(1)
 \longrightarrow T_p(B)\otimes C
 \xrightarrow{\alpha_B} \omega_{B^D}\otimes C
 \longrightarrow0 ,
\]
where the first arrow is the transpose of \(\alpha_{B^D}\) via
Cartier duality. In particular the kernel of \(\alpha_B\) is the
annihilator of the kernel of \(\alpha_{B^D}\) under the Tate pairing.
Naturality gives
\(\alpha_{B^D}\circ T_p(\lambda)=(\lambda^D)^*\circ\alpha_B\).
Thus \(\ker\alpha_B\subseteq\ker H\). Both kernels have dimension
\(g\), since \(T_p(\lambda)\) is an isomorphism and both Hodge--Tate
maps are surjective over \(C\) with rank \(g\). They coincide.
Using the alternating pairing given by \(\lambda\), the annihilator
identity becomes \((\ker H)^\perp=\ker H\).
This is the required orthogonality, including at \(p=2\).
Only rational exactness of the displayed complex is used here;
for the integral cokernel bound we retain (HS8) and its factor \(4\).
\end{proof}

\begin{theorem}[An exact flag sequence modulo the original loss]
\label{thm:exact-corrected-flag-sequence}
Keep the original correction prime \(p\), \(\mathrm{HTFree}\) model
and the original flag with \(\mathrm{RootLines}\). On its
\(p^{R-2}\)-thickened fibre, for each CM-type embedding \(\sigma\)
there is an exact sequence of vector bundles
\[
 0\longrightarrow P_\sigma
 \longrightarrow \mathcal O^{N}
 \longrightarrow Q_\sigma
 \longrightarrow0,
 \qquad
 P_\sigma=\sigma(\mathfrak D_E)^{-1}
                  Q_{\bar\sigma}^{\vee},\quad N=p_{\rm sig}+q_{\rm sig}.
 \tag{HF1}
\]
Here \(Q_\sigma\) denotes the restricted image of the original
corrected cotangent module, with the \emph{conjugate} coefficient
action on differentials; all bundles in (HF1) are reduced modulo
\(p^{R-2}\). Thus \(P_\sigma\) is the original selected corrected
tangent bundle and \(Q_\sigma\) is the dual of the other corrected
tangent bundle. The sequence uses the original full-level basis
and normalized Weil pairing. It supplies \(\mathrm{TrivFlagOnHodge}\)
and the correction-prime \(\mathrm{ModelNefStatement}\) on the
actual model.
\end{theorem}

\begin{proof}
Use the convention on the action on differentials just stated, so
the names \(Q_\sigma\) here correspond to the source
\(\mathrm{omHT}\), not to an un-conjugated dual of the
\(\sigma\)-eigenspace. Before reduction work on the original model
over \(\mathbb Z_{(p)}\). By (HS8), and the inclusions proved on
the locally free charts,
\[
 p^R\omega_G\ \subset\ p^{R-2}\mathcal N_p,\qquad
 \mathcal N_p=\langle h(t)\rangle+p^R\omega_G ,
 \tag{HF2}
\]
where \(h(t)\) are local lifts of the actual character differentials.
For (HF2), \(\kappa_p\mid p^2\) and
\(\kappa_p\omega_G\subseteq\mathcal N_p\).
The containment is also valid with the original projected image in
place of \(\mathcal N_p\). In particular all the character
differentials are canonically defined in
\(\mathcal N_p/p^m\mathcal N_p\), irrespective of the lifts.
Their additivity and coefficient linearity hold there: both were
identities modulo \(p^R\omega_G\) before using (HF2).

Choose the original \(N\) generators of the free
\(\mathcal O_E/p^R\mathcal O_E\)-module of generic torsion.
Their projected differentials now give a global surjective map
\[
 q_\sigma:(\mathcal O/p^m)^N\longrightarrow Q_\sigma/p^m Q_\sigma .
 \tag{HF3}
\]
Indeed after projection to the \(\sigma\)-eigenspace the coefficients
in \(\mathcal O_E\) act through \(\sigma\), and the additional
\(p^R\omega_G\) generators disappear by (HF2).
All embeddings of integers are integral on the normal charts.
This argument does not assume that the idempotent projections are
integral on the uncorrected cotangent module. It uses precisely
their locally free images.

Here is the integral normalization of the adjoint. Put
\(D_p=\mathcal O_E\otimes\mathbb Z_{(p)}\).
Its inverse different is the trace dual of \(D_p\), and it is
principal because \(D_p\) is semilocal Dedekind. Choose once a
generator \(d\) of its different. The original perfect alternating
balanced pairing of the full level, after choosing the original
primitive \(p^R\)-th root of unity, is the trace of a unique
sesquilinear pairing with values in
\(\mathfrak D_E^{-1}D_p/p^R\mathfrak D_E^{-1}D_p\).
Indeed the trace-dual identification applied in one variable
constructs that pairing; the balanced identity gives its conjugate
linearity in the other variable. The perfectness of the trace
pairing over \(\mathbb Z/p^R\) implies perfectness of this
sesquilinear pairing over \(D_p/p^R\).
Multiplication by \(d\) makes its matrix invertible over that
ring. Evaluation at \(\sigma\) therefore gives an invertible
bilinear matrix pairing the two conjugate labelled copies of
\((\mathcal O/p^m)^N\).
This is the different-normalized Gram matrix of
\(\mathrm{HTChart.gram}\); it is defined by the constant full-level
pairing, not by choosing a fresh Gram form on each chart.

The dual of \(q_{\bar\sigma}\) composed with the inverse of this
matrix defines an injection \(i_\sigma:P_\sigma/p^m\to
(\mathcal O/p^m)^N\). In this formula a member of
\(P_\sigma=\sigma(d)^{-1}Q_{\bar\sigma}^{\vee}\) is first multiplied
by \(\sigma(d)\) to give a member of the ordinary dual, as required
by the different-normalized Gram. Equivalently the defining
adjoint equality is the one for the unscaled pairing and evaluation
of the member of \(P_\sigma\). This construction is independent of
the choice of different generator. Since \(Q_{\bar\sigma}\) is
locally free and \(q_{\bar\sigma}\) is surjective, \(i_\sigma\) is
locally a split injection, including over a nonreduced thickening.

We claim \(q_\sigma i_\sigma=0\). It is enough to check this at
the height-one primes above \(p\) of the original normal chart.
Indeed locally lift the composite to a map of projective modules
before reduction; its image in the fraction field is divisible
by \(p^m\) in the target if and only if its quotient by \(p^m\)
belongs to the target at every height-one prime. This is the
intersection property of reflexive modules over a normal
noetherian domain. At height-one primes not over \(p\) the
divisibility condition is automatic.

At a height-one prime above \(p\), extend faithfully flatly from
the DVR to the ring of the complete algebraically closed valued
field. Construct \(B\) by Lemma~\ref{lem:corrected-ht-test-tower}.
Lift the original level generators to an
\(\mathcal O_E\otimes\mathbb Z_p\)-basis of \(T_p(B)\);
Nakayama's lemma, using their basis property modulo \(p^R\),
justifies the lift. Lift also the chosen root of unity compatibly
to all levels. The Tate pairing reduces to the original one
modulo \(p^R\); its trace-dual pairing and its normalized Gram
do so as well.

The image of the actual Tate Hodge--Tate map
\(\alpha_{B^D}\circ T_p(\lambda)\) over \(O_C\) is exactly
\(\mathcal N_p\otimes O_C\). Its reduction in \(\omega_G/p^R\)
is the original character differential map, since at each level
the identity copy of \(B_n\) is \(\mathcal H_n\).
Its cokernel is annihilated by \(\kappa_p\): apply the integral
Hodge--Tate cokernel bound (or the compatible finite bounds and
compact finite generation of the image) to \(B^D\).
In particular its image contains \(p^R\omega_G\).
The equality of the two images follows.
The finite free projected modules commute with this flat
base change, as images of a fixed rational linear map.

Over \(O_C\) therefore one can use the actual Tate Hodge--Tate
map and actual Tate Gram matrix before reduction. Their maps
on the projected image lattices are integral and give exactly
(HF3) and its adjoint after reduction modulo \(p^m\).
The kernel of the rational Hodge--Tate map is Lagrangian by
Lemma~\ref{lem:corrected-ht-test-tower}, so the adjoint of the
conjugate labelled quotient has image in the kernel of the
labelled quotient. The different normalization merely scales
that adjoint by a nonzero constant over \(C\).
The composite of the actual integral maps is consequently zero
over \(C\), hence over \(O_C\). The congruence between actual
Tate differentials and the chosen level differentials modulo
\(p^R\omega_G\), together with (HF2), proves zero composition
for the original reduced maps. Faithful flatness and the
reflexive-module argument prove the claim on the original chart.

The ranks of \(P_\sigma\) and \(Q_\sigma\) are respectively
\(p_{\rm sig}\) and \(q_{\rm sig}\), so they sum to \(N\).
The split injection and surjection just proved are exact:
locally quotient the free middle bundle by the split first
bundle. The induced surjection to \(Q_\sigma\) is between free
modules of the same rank, hence is an isomorphism over the
possibly nonreduced local ring. This proves (HF1).
The maps glue because the original torsion labels, Weil pairing,
different ideal and classes in (HF2) are global.

Now use the actual flag in \(P_\sigma\), and in \(Q_\sigma\)
use the reversed dual flag of the other corrected tangent bundle.
The successive inverse images of this second flag under (HF1),
following the first flag, give a full flag of subbundles of the
trivial middle bundle modulo \(p^m\). Its graded line bundles are
exactly the original root lines, including the inverses and reversal
in \(\mathrm{rootIdx},\mathrm{rootSgn}\). The flag restricts to every
special-fibre component without any new base-field extension.
In particular it supplies the source's \(\mathrm{TrivFlagAt}\)
on those components. In the coordinate model one can choose a flag
basis on each original affine chart after localizing as needed;
the prescribed root-line generators fix the diagonal transition
factors. The change matrices are upper triangular modulo \(p^m\),
and the original root monomial ratios are their diagonal entries.
Equivalently the exact sequence of subbundles is the geometric
form of the original trivial-flag assertion; no change to the
coordinate model or to its degree is necessary.

For a complete integral curve on a special-fibre component,
write \(L_i\) for the successive quotients of this full flag and
\(F_i\) for its \(i\)-th subbundle (\(F_0=0,F_N=\mathcal O^N\)).
The determinant of \(F_i\) injects into a trivial exterior-power
bundle, so \(\deg F_i\le0\): a nonzero component of the induced
map to a trivial line gives an effective divisor after normalization
of the curve. Also \(\deg F_N=0\). Thus for decreasing real
weights \(\mu_1\ge\cdots\ge\mu_N\), summation by parts gives
\[
 -\sum_{i=1}^N\mu_i\deg L_i
 =-\sum_{i=1}^{N-1}(\mu_i-\mu_{i+1})\deg F_i\ \ge0 .
 \tag{HF4}
\]
This is the sign convention of
\(\mathrm{RootModel.coef}(\mu,0)=-\mu\).
Sum over the original CM-type labels. It proves precisely
\(\mathrm{ModelNefStatement}\) at the correction primes.
The Whitney formula also gives
\(\prod_i(1+c_1(L_i))=1\) on each component, and the positive-degree
power sums of these roots vanish there by the Newton identities.
\end{proof}

The proof establishes the exact Lagrangian comparison and the
correction-prime general flag positivity, rather than only an
annihilator inclusion. It does not supply the distinct \(\mathrm{M2}\)
condition away from the correction set, or the Frobenius-linked
vertical correction identities comparing two levels. Those remain
necessary for the complete weighted mean. No numerical change to
the base-test bound of Section~\ref{sec:geometry-height-assembly}
is asserted here.
